\begin{block}{Translation-Based Annotation}
Uses Kaneko's modern Japanese gloss as interpretive evidence. A locally run LLM
(\texttt{qwen2.5} / Ollama) first aligns the poem's five segments with the modern
Japanese translation, then finds sentence- or discourse-level breaks in the
translation, and finally maps them back onto segment positions in the original
poem. This method emphasizes ``the places read as broken in interpretive
practice.''

Human review and TEI output: automatically generated candidates are first
written to a per-poem cache file, then checked, added to, or corrected by a
human reviewer. The confirmed result is written back to TEI/XML, so the data
stays traceable to the evidence behind each kugire judgment.

\begin{figure}[H]\centering
\includesvg[inkscapelatex=false, width=0.9\textwidth]{../images/fig-annotation-pipeline-tikz}
\setcounter{figure}{5}
\caption{Annotation pipeline from poem and translation, through LLM
suggestion and human review, to TEI/XML output.}
\end{figure}
\end{block}
